% GENERATED FILE. Edit canonical lessons, then run npm run generate:books.
% canonical-source-hash: fnv1a64:c8eba9d54eca98b2
% canonical-lessons: ML-C526-kulannuka, ML-C526-ksobhikkuka, ML-C526-korittarikkuka, ML-C526-rasikkuka, ML-C526-musiyuka

\chapter{Stumped, Stirred and Amused}
\label{ch:ml-stumped-stirred-and-amused}
\hlchaptermodality{526}

\begin{chapteropening}
\textbf{By the end of this chapter:} \emph{I can say to be stumped, to get agitated, to thrill, to be amused and to grow weary.}

Complete the last lesson of chapter 526: I can say to be stumped, to get agitated, to thrill, to be amused and to grow weary.
\end{chapteropening}

\section[\texorpdfstring{\ml{കുഴങ്ങുക} (kuḻaṅṅuka)}{കുഴങ്ങുക (kuḻaṅṅuka)}]{\ml{കുഴങ്ങുക} (kuḻaṅṅuka) — to be stumped; to be at a loss}

\label{lesson:ML-C526-kulannuka}



\begin{quote}
Before the new one: say the Malayalam for to burst out laughing, then the Malayalam for confusion.
\end{quote}

\subsection*{You'll want to know: \ml{കുഴങ്ങുക}}
\textbf{\ml{കുഴങ്ങുക}} — \emph{kuḻaṅṅuka} — \textquotedblleft{}to be stumped; to be at a loss\textquotedblright{}.

\begin{quote}
\ml{ചോദ്യം} \ml{കേട്ട്} \ml{കുട്ടി} \ml{കുഴങ്ങി}.
\end{quote}

\begin{quote}
\emph{cōdyaṁ kēṭṭŭ kuṭṭi kuḻaṅṅi.}
\end{quote}

\begin{quote}
The child was stumped by the question.
\end{quote}

Past: \ml{കുഴങ്ങി}, \emph{kuḻaṅṅi}.

Often \ml{കുഴങ്ങിപ്പോയി}, was left at a loss. \ml{ആശയക്കുഴപ്പം} shares its root.

\begin{cousinweb}
\textbf{\ml{കുഴങ്ങുക}} (\emph{kuḻaṅṅuka}) is native Dravidian, the verb beside \ml{കുഴപ്പം}, \textquotedblleft{}intricacy, trouble\textquotedblright{} (DEDR 1817).
\end{cousinweb}

\begin{grammarlens}[title={What you've built}]
One more word for describing what you lived and felt.
\end{grammarlens}

\subsection*{Your turn}
\begin{itemize}
\raggedright
  \item \emph{Say it:} \emph{kuḻaṅṅuka}
  \item \emph{Say it:} \emph{kuḻaṅṅuka}, once more
  \item \emph{Say it:} \emph{cōdyaṁ kēṭṭŭ kuṭṭi kuḻaṅṅi.} — The child was stumped by the question.
  \item \emph{Recall:} say the Malayalam for to burst out laughing, then the Malayalam for confusion, then say \emph{kuḻaṅṅuka} again
\end{itemize}

\begin{tcolorbox}[breakable,colback=teal!4,colframe=teal!35!black,title={Before you move on}]
What does \ml{കുഴങ്ങുക} mean? (\textquotedblleft{}To be stumped; to be at a loss\textquotedblright{}.) Say it once more.
\end{tcolorbox}
\section[\texorpdfstring{\ml{ക്ഷോഭിക്കുക} (kṣōbhikkuka)}{ക്ഷോഭിക്കുക (kṣōbhikkuka)}]{\ml{ക്ഷോഭിക്കുക} (kṣōbhikkuka) — to get agitated; to flare up}

\label{lesson:ML-C526-ksobhikkuka}



\begin{quote}
Before the new one: say the Malayalam for confusion, then the Malayalam for to be stumped.
\end{quote}

\subsection*{You'll want to know: \ml{ക്ഷോഭിക്കുക}}
\textbf{\ml{ക്ഷോഭിക്കുക}} — \emph{kṣōbhikkuka} — \textquotedblleft{}to get agitated; to flare up\textquotedblright{}.

\begin{quote}
\ml{ശമ്പളം} \ml{കിട്ടാതെ} \ml{വന്നപ്പോൾ} \ml{തൊഴിലാളികൾ} \ml{ക്ഷോഭിച്ചു}.
\end{quote}

\begin{quote}
\emph{śampaḷaṁ kiṭṭāte vannappōḷ toḻilāḷikaḷ kṣōbhiccu.}
\end{quote}

\begin{quote}
When their wages did not come, the workers grew agitated.
\end{quote}

Past: \ml{ക്ഷോഭിച്ചു}, \emph{kṣōbhiccu}.

Formal; the noun is \ml{ക്ഷോഭം}, agitation, also of the sea.

\begin{cousinweb}
\textbf{\ml{ക്ഷോഭിക്കുക}} (\emph{kṣōbhikkuka}) is built from Sanskrit \emph{kṣobha}, \textquotedblleft{}shaking, agitation, emotion\textquotedblright{} (Monier-Williams), with the Malayalam verb ending -\ml{ിക്കുക}.
\end{cousinweb}

\begin{grammarlens}[title={What you've built}]
One more word for describing what you lived and felt.
\end{grammarlens}

\subsection*{Your turn}
\begin{itemize}
\raggedright
  \item \emph{Say it:} \emph{kṣōbhikkuka}
  \item \emph{Say it:} \emph{kṣōbhikkuka}, once more
  \item \emph{Say it:} \emph{śampaḷaṁ kiṭṭāte vannappōḷ toḻilāḷikaḷ kṣōbhiccu.} — When their wages did not come, the workers grew agitated.
  \item \emph{Recall:} say the Malayalam for confusion, then the Malayalam for to be stumped, then say \emph{kṣōbhikkuka} again
\end{itemize}

\begin{tcolorbox}[breakable,colback=teal!4,colframe=teal!35!black,title={Before you move on}]
What does \ml{ക്ഷോഭിക്കുക} mean? (\textquotedblleft{}To get agitated; to flare up\textquotedblright{}.) Say it once more.
\end{tcolorbox}
\section[\texorpdfstring{\ml{കോരിത്തരിക്കുക} (kōrittarikkuka)}{കോരിത്തരിക്കുക (kōrittarikkuka)}]{\ml{കോരിത്തരിക്കുക} (kōrittarikkuka) — to thrill; to get goosebumps}

\label{lesson:ML-C526-korittarikkuka}



\begin{quote}
Before the new one: say the Malayalam for to be stumped, then the Malayalam for to get agitated.
\end{quote}

\subsection*{You'll want to know: \ml{കോരിത്തരിക്കുക}}
\textbf{\ml{കോരിത്തരിക്കുക}} — \emph{kōrittarikkuka} — \textquotedblleft{}to thrill; to get goosebumps\textquotedblright{}.

\begin{quote}
\ml{ആ} \ml{പാട്ട്} \ml{കേട്ട്} \ml{ഞാൻ} \ml{കോരിത്തരിച്ചു}.
\end{quote}

\begin{quote}
\emph{ā pāṭṭŭ kēṭṭŭ ñān kōrittariccu.}
\end{quote}

\begin{quote}
I got goosebumps listening to that song.
\end{quote}

Past: \ml{കോരിത്തരിച്ചു}, \emph{kōrittariccu}.

The thrill one feels on the skin, as when the hair stands on end.

\begin{cousinweb}
\textbf{\ml{കോരിത്തരിക്കുക}} (\emph{kōrittarikkuka}) looks like a native compound verb, but DEDR has no entry that clearly holds it, so its origin is not settled.
\end{cousinweb}

\begin{grammarlens}[title={What you've built}]
One more word for describing what you lived and felt.
\end{grammarlens}

\subsection*{Your turn}
\begin{itemize}
\raggedright
  \item \emph{Say it:} \emph{kōrittarikkuka}
  \item \emph{Say it:} \emph{kōrittarikkuka}, once more
  \item \emph{Say it:} \emph{ā pāṭṭŭ kēṭṭŭ ñān kōrittariccu.} — I got goosebumps listening to that song.
  \item \emph{Recall:} say the Malayalam for to be stumped, then the Malayalam for to get agitated, then say \emph{kōrittarikkuka} again
\end{itemize}

\begin{tcolorbox}[breakable,colback=teal!4,colframe=teal!35!black,title={Before you move on}]
What does \ml{കോരിത്തരിക്കുക} mean? (\textquotedblleft{}To thrill; to get goosebumps\textquotedblright{}.) Say it once more.
\end{tcolorbox}
\section[\texorpdfstring{\ml{രസിക്കുക} (rasikkuka)}{രസിക്കുക (rasikkuka)}]{\ml{രസിക്കുക} (rasikkuka) — to be amused; to have fun}

\label{lesson:ML-C526-rasikkuka}



\begin{quote}
Before the new one: say the Malayalam for to get agitated, then the Malayalam for to thrill.
\end{quote}

\subsection*{You'll want to know: \ml{രസിക്കുക}}
\textbf{\ml{രസിക്കുക}} — \emph{rasikkuka} — \textquotedblleft{}to be amused; to have fun\textquotedblright{}.

\begin{quote}
\ml{കുട്ടികൾ} \ml{നാടകം} \ml{കണ്ട്} \ml{ഒരുപാട്} \ml{രസിച്ചു}.
\end{quote}

\begin{quote}
\emph{kuṭṭikaḷ nāṭakaṁ kaṇṭŭ orupāṭŭ rasiccu.}
\end{quote}

\begin{quote}
The children had great fun watching the play.
\end{quote}

Past: \ml{രസിച്ചു}, \emph{rasiccu}.

Its noun \ml{രസം}, fun, also gives \ml{രസകരമായ}, interesting.

\begin{cousinweb}
\textbf{\ml{രസിക്കുക}} (\emph{rasikkuka}) is built from Sanskrit \emph{rasa}, \textquotedblleft{}sap, taste, relish\textquotedblright{} (CDIAL 10650), with the Malayalam verb ending -\ml{ിക്കുക}.
\end{cousinweb}

\begin{grammarlens}[title={What you've built}]
One more word for describing what you lived and felt.
\end{grammarlens}

\subsection*{Your turn}
\begin{itemize}
\raggedright
  \item \emph{Say it:} \emph{rasikkuka}
  \item \emph{Say it:} \emph{rasikkuka}, once more
  \item \emph{Say it:} \emph{kuṭṭikaḷ nāṭakaṁ kaṇṭŭ orupāṭŭ rasiccu.} — The children had great fun watching the play.
  \item \emph{Recall:} say the Malayalam for to get agitated, then the Malayalam for to thrill, then say \emph{rasikkuka} again
\end{itemize}

\begin{tcolorbox}[breakable,colback=teal!4,colframe=teal!35!black,title={Before you move on}]
What does \ml{രസിക്കുക} mean? (\textquotedblleft{}To be amused; to have fun\textquotedblright{}.) Say it once more.
\end{tcolorbox}
\section[\texorpdfstring{\ml{മുഷിയുക} (muṣiyuka)}{മുഷിയുക (muṣiyuka)}]{\ml{മുഷിയുക} (muṣiyuka) — to grow weary; to get grubby}

\label{lesson:ML-C526-musiyuka}



\begin{quote}
Before the new one: say the Malayalam for to thrill, then the Malayalam for to be amused.
\end{quote}

\subsection*{You'll want to know: \ml{മുഷിയുക}}
\textbf{\ml{മുഷിയുക}} — \emph{muṣiyuka} — \textquotedblleft{}to grow weary; to get grubby\textquotedblright{}.

\begin{quote}
\ml{ഒരുപാട്} \ml{നേരം} \ml{കാത്തിരുന്ന്} \ml{ഞങ്ങൾ} \ml{മുഷിഞ്ഞു}.
\end{quote}

\begin{quote}
\emph{orupāṭŭ nēraṁ kāttirunnŭ ñaṅṅaḷ muṣiññu.}
\end{quote}

\begin{quote}
We grew weary of waiting so long.
\end{quote}

Past: \ml{മുഷിഞ്ഞു}, \emph{muṣiññu}.

Said of clothes it means to get grubby: \ml{മുഷിഞ്ഞ} \ml{ഷർട്ട്}, a grubby shirt. The past is \ml{മുഷിഞ്ഞു}.

\begin{cousinweb}
\textbf{\ml{മുഷിയുക}} (\emph{muṣiyuka}) is native Dravidian: DEDR 4903 gives the Malayalam verb with both senses, to grow weary and to get dirty (DEDR 4903).
\end{cousinweb}

\begin{grammarlens}[title={What you've built}]
One more word for describing what you lived and felt.
\end{grammarlens}

\subsection*{Your turn}
\begin{itemize}
\raggedright
  \item \emph{Say it:} \emph{muṣiyuka}
  \item \emph{Say it:} \emph{muṣiyuka}, once more
  \item \emph{Say it:} \emph{orupāṭŭ nēraṁ kāttirunnŭ ñaṅṅaḷ muṣiññu.} — We grew weary of waiting so long.
  \item \emph{Recall:} say the Malayalam for to thrill, then the Malayalam for to be amused, then say \emph{muṣiyuka} again
\end{itemize}

\begin{tcolorbox}[breakable,colback=teal!4,colframe=teal!35!black,title={Before you move on}]
What does \ml{മുഷിയുക} mean? (\textquotedblleft{}To grow weary; to get grubby\textquotedblright{}.) Say it once more.
\end{tcolorbox}
